\section{Discussion}
\label{sec:discussion}

The seven-system results establish that terminal decision procedures can improve a
range of medical question-answering systems under a common native evaluation contract.
The benefit is not confined to the two MedRAG configurations: direct answering,
MedRGAG, MA-RAG, MedGENIE, and TC-RAG also improve throughout the three declared
development strata. Preserving these results alongside their controls is necessary to
understand the research outcome. Mixed contributions from calibration, or incomplete
external validation, do not erase the measured gain of the complete terminal head.

What changes depends on the task. The diagnostic component provides an explicit answer
space and a new original-model decision; the rule component executes support conditions
and combines them with the original answer. These are distinct interventions, and R1's
reuse of R2/R3 components demonstrates useful transfer across support requirements
without constituting three independent algorithms. A common component can also preserve
method differences through the original context, UNKNOWN inheritance, and fallback.
The evidence supports system improvement with explicit attribution limits, rather than
a claim that all failures share one cause.

The later CF-MoA core combines the catalog, rules, and repeated readout in two
configurations. Its improvement concerns that broader composition and its additional
computation. It neither replaces the seven-system findings nor extends the full CF-MoA
result to five untested configurations. The composition's gain also cannot be reassigned
to candidate verification introduced afterward.

Additional candidate verification poses a distinct selection problem. By design, C3 cannot
recover an answer key absent from its pool, and it can replace a correct selected answer
with an incorrect existing candidate. The repair and damage counts establish that
successful and harmful switches both occur; they do not measure the frequency of missing
correct candidates or diagnose the medical cause of each mistake. Pool quality and
selector quality should therefore be measured separately. Enlarging the pool is also not
automatically beneficial: the C4 control changes both candidate opportunity and failure
exposure, with its unavailable outputs retained in the result.

Rationales provide extra information to the verifier, but they can contain unsupported
statements or persuasive explanations for a wrong key. Their positive incremental point
estimates over C2 justify further study of how they affect selection. They do not certify
rationale truthfulness, establish calibrated confidence, or guarantee improvement over
majority-based readout. A direct evaluation of rationale validity and an independent
replication of pool generation would be needed to test those stronger interpretations.

The source and pair analyses constrain the use of an aggregate gain. A positive overall
score can coexist with harm on a particular source, and a changed overall mean can coexist
with unchanged respond-pair performance. No arm was selected separately for each source or
backbone after inspecting these outcomes. Keeping the adopted core as the common default
preserves the distinction between an observed local gain and a deployment or
generalization decision.

Finally, the five-role description is useful only when tied to actual computation. Shared
rule addition and removal are two functions of one execution. A model-based executor with
zero coverage cannot explain a clinical benchmark gain. Recording these facts makes the
architecture interpretable without implying that every question benefits from five-agent
collaboration.
